\documentclass[../../../include/open-logic-section]{subfiles}

\begin{document}

\olfileid{int}{tab}{prf}

% Part: intuitionistic-logic
如 \olref{int}{tab}{mon} 中所述，直觉主义 克里普克 模型中的满足关系是单调的：如果 $\mathfrak{M}, \Gamma \models !A$ 且 $\Gamma R \Delta$，则 $\mathfrak{M}, \Delta \models !A$。我们现在通过对~$!A$ 的结构归纳来证明这一点。

\paragraph{归纳基始：}如果 $!A$ 是原子公式，则由赋值的单调性条件可知结论立即成立。如果 $!A \ident \lfalse$，则 $\mathfrak{M},\Gamma \not\models \lfalse$（根据定义），因此“如果 $\mathfrak{M}, \Gamma \models \lfalse$ 且 $\Gamma R \Delta$，则 $\mathfrak{M}, \Delta \models \lfalse$”的前件为假，从而该条件句平凡地为真。

\paragraph{归纳步骤：}我们考虑 $!A$ 的各种可能形式。
\begin{enumerate}
\item 如果 $!A \ident !B \land !C$，且 $\mathfrak{M}, \Gamma \models !A$，则 $\mathfrak{M}, \Gamma \models !B$ 且 $\mathfrak{M}, \Gamma \models !C$。根据归纳假设，如果 $\Gamma R \Delta$，则 $\mathfrak{M}, \Delta \models !B$ 且 $\mathfrak{M}, \Delta \models !C$，因此 $\mathfrak{M}, \Delta \models !B \land !C$。

\item 如果 $!A \ident !B \lor !C$，且 $\mathfrak{M}, \Gamma \models !A$，则 $\mathfrak{M}, \Gamma \models !B$ 或 $\mathfrak{M}, \Gamma \models !C$。根据归纳假设，如果 $\Gamma R \Delta$，则 $\mathfrak{M}, \Delta \models !B$ 或 $\mathfrak{M}, \Delta \models !C$，因此 $\mathfrak{M}, \Delta \models !B \lor !C$。

\item 如果 $!A \ident !B \lif !C$，且 $\mathfrak{M}, \Gamma \models !A$，则对所有满足 $\Gamma R \Delta$ 的 $\Delta$，如果 $\mathfrak{M}, \Delta \models !B$，则 $\mathfrak{M}, \Delta \models !C$。现在假设 $\Gamma R \Delta'$。为了证明 $\mathfrak{M}, \Delta' \models !B \lif !C$，任取 $\Delta''$ 满足 $\Delta' R \Delta''$ 且 $\mathfrak{M}, \Delta'' \models !B$。由于 $R$ 是传递的，有 $\Gamma R \Delta''$。由假设，因 $\Gamma R \Delta''$ 且 $\mathfrak{M}, \Delta'' \models !B$，可得 $\mathfrak{M}, \Delta'' \models !C$。所以 $\mathfrak{M}, \Delta' \models !B \lif !C$。

\item 如果 $!A \ident \lforall[x][!B]$，且 $\mathfrak{M}, \Gamma \models !A$，则对所有满足 $\Gamma R \Delta$ 的 $\Delta$ 及所有 $d \in D(\Delta)$，有 $\mathfrak{M}, \Delta \models !B[x\backslash d]$。现在假设 $\Gamma R \Delta'$。为了证明 $\mathfrak{M}, \Delta' \models \lforall[x][!B]$，任取 $\Delta''$ 满足 $\Delta' R \Delta''$ 及 $d \in D(\Delta'')$。由于 $R$ 是传递的，有 $\Gamma R \Delta''$。由假设，可得 $\mathfrak{M}, \Delta'' \models !B[x\backslash d]$。所以 $\mathfrak{M}, \Delta' \models \lforall[x][!B]$。

\item 如果 $!A \ident \lexists[x][!B]$，且 $\mathfrak{M}, \Gamma \models !A$，则存在某个 $d \in D(\Gamma)$ 使得 $\mathfrak{M}, \Gamma \models !B[x\backslash d]$。根据归纳假设，如果 $\Gamma R \Delta$，则 $\mathfrak{M}, \Delta \models !B[x\backslash d]$。由于 $d \in D(\Gamma) \subseteq D(\Delta)$，因此 $\mathfrak{M}, \Delta \models \lexists[x][!B]$。
\end{enumerate}
在合取、析取与存在量词的情况中，我们都应用了归纳假设（分别应用于 $!B$、$!C$，或存在情况下的 $!B[x\backslash d]$）。这就完成了证明。

% Chapter: tableaux
% Section: proofs

\olsection{直觉主义逻辑的\usetoken{P}{tableau}}

\begin{ex}
  我们给出一个封闭的分析表来证明 $(!A \land !B) \lif !C \Proves 
  !A \lif (!B \lif !C)$。
  \begin{oltableau}
    [\pFmla{\True}{(\formula{A} \land \formula{B}) \lif \formula{C}}{1},
      just =\TAss
      [\pFmla{\False}{\formula{A} \lif (\formula{B} \lif \formula{C})}{1},
        just =\TAss
        [\pFmla{\True}{\formula{A}}{1.1},
          just = {\TRule{\False}{\lif}[2]}
          [\pFmla{\False}{\formula{B} \lif \formula{C}}{1.1},
            just = {\TRule{\False}{\lif}[2]}
            [\pFmla{\True}{\formula{B}}{1.1.1},
              just = {\TRule{\False}{\lif}[4]}
              [\pFmla{\False}{\formula{C}}{1.1.1},
                just = {\TRule{\False}{\lif}[4]}
                [\pFmla{\False}{\formula{A} \land \formula{B}}{1.1.1},
                  just = {\TRule{\True}{\lif}[1]}
                  [\pFmla{\False}{\formula{A}}{1.1.1},
                    just= {\TRule{\False}{\land}[4]}, close]
                  [\pFmla{\False}{\formula{B}}{1.1.1},
                    just= {\TRule{\False}{\land}[4]}, close]]
                [\pFmla{\True}{\formula{C}}{1.1.1},
                  just = {\TRule{\True}{\lif}[1]}, close]]
            ]
          ]
        ]
      ]
    ]
  \end{oltableau}
\end{ex}

\begin{prob}
  构造封闭的直觉主义分析表以证明下列公式：
  \begin{enumerate}
    \item $\Proves !A \lif (!B \lif !A)$
    \item $\Proves \lnot(!A \land \lnot !A)$
    \item $!A \lif (!B \lif !C) \Proves (!A \land !B) \lif !C$
    \item $\lnot !A \lor \lnot !B \Proves \lnot(!A \land !B)$
  \end{enumerate}
\end{prob}

\end{document}